\documentclass[10pt,a4paper]{amsart}
\usepackage{amsmath,amssymb,amsthm,graphicx}
% Retained minimal preamble for the Stacks Project chapter more-algebra.tex (More on Algebra).
% The chapter file pinned-original-source.tex carries only \input{preamble} on its first line, so no
% preamble text lives inside the chapter itself.  The Stacks Project uses one single shared preamble for
% every chapter; the snapshot stages it identically as sources/stacks-chow/preamble.tex and
% sources/stacks-etale-cohomology/preamble.tex.  The block below is copied mechanically, line for line,
% from sources/stacks-chow/preamble.tex lines 14-18 (the Xy-pic setup) and 154-217 (the native theorem
% declarations together with the mathematical macros).  Nothing was added, removed or redefined; the
% wrapper class and the amsmath/amssymb/amsthm/graphicx base packages are supplied by the assembler.
% Chapter: more-algebra.tex -- More on Algebra.
% For commutative diagrams we use Xy-pic
\usepackage[all]{xy}

% We use 2cell for 2-commutative diagrams.
\xyoption{2cell}

% Theorem environments.
%
\theoremstyle{plain}
\newtheorem{theorem}[subsection]{Theorem}
\newtheorem{proposition}[subsection]{Proposition}
\newtheorem{lemma}[subsection]{Lemma}

\theoremstyle{definition}
\newtheorem{definition}[subsection]{Definition}
\newtheorem{example}[subsection]{Example}
\newtheorem{exercise}[subsection]{Exercise}
\newtheorem{situation}[subsection]{Situation}

\theoremstyle{remark}
\newtheorem{remark}[subsection]{Remark}
\newtheorem{remarks}[subsection]{Remarks}

\numberwithin{equation}{subsection}

% Macros
%
\def\lim{\mathop{\mathrm{lim}}\nolimits}
\def\colim{\mathop{\mathrm{colim}}\nolimits}
\def\Spec{\mathop{\mathrm{Spec}}}
\def\Hom{\mathop{\mathrm{Hom}}\nolimits}
\def\Ext{\mathop{\mathrm{Ext}}\nolimits}
\def\SheafHom{\mathop{\mathcal{H}\!\mathit{om}}\nolimits}
\def\SheafExt{\mathop{\mathcal{E}\!\mathit{xt}}\nolimits}
\def\Sch{\mathit{Sch}}
\def\Mor{\mathop{\mathrm{Mor}}\nolimits}
\def\Ob{\mathop{\mathrm{Ob}}\nolimits}
\def\Sh{\mathop{\mathit{Sh}}\nolimits}
\def\NL{\mathop{N\!L}\nolimits}
\def\CH{\mathop{\mathrm{CH}}\nolimits}
\def\proetale{{pro\text{-}\acute{e}tale}}
\def\etale{{\acute{e}tale}}
\def\QCoh{\mathit{QCoh}}
\def\Ker{\mathop{\mathrm{Ker}}}
\def\Im{\mathop{\mathrm{Im}}}
\def\Coker{\mathop{\mathrm{Coker}}}
\def\Coim{\mathop{\mathrm{Coim}}}

% Boxtimes
%
\DeclareMathSymbol{\boxtimes}{\mathbin}{AMSa}{"02}

%
% Macros for moduli stacks/spaces
%
\def\QCohstack{\mathcal{QC}\!\mathit{oh}}
\def\Cohstack{\mathcal{C}\!\mathit{oh}}
\def\Spacesstack{\mathcal{S}\!\mathit{paces}}
\def\Quotfunctor{\mathrm{Quot}}
\def\Hilbfunctor{\mathrm{Hilb}}
\def\Curvesstack{\mathcal{C}\!\mathit{urves}}
\def\Polarizedstack{\mathcal{P}\!\mathit{olarized}}
\def\Complexesstack{\mathcal{C}\!\mathit{omplexes}}
% \Pic is the operator that assigns to X its picard group, usage \Pic(X)
% \Picardstack_{X/B} denotes the Picard stack of X over B
% \Picardfunctor_{X/B} denotes the Picard functor of X over B
\def\Pic{\mathop{\mathrm{Pic}}\nolimits}
\def\Picardstack{\mathcal{P}\!\mathit{ic}}
\def\Picardfunctor{\mathrm{Pic}}
\def\Deformationcategory{\mathcal{D}\!\mathit{ef}}
\usepackage{cleveref}
\crefname{originallemma}{original lemma}{original lemmas}
\Crefname{originallemma}{Original lemma}{Original lemmas}
\makeatletter
\newlabel{lemma-sequence-powers-pro-bounded}{{0G3M}{source}{}{sourcepointer-1}{}}
\newlabel{lemma-sequence-powers-pro-bounded@cref}{{[originallemma][0][]0G3M}{[1][0][]source}{}{}{}}
\makeatother
\crefname{originallemma}{original lemma}{original lemmas}
\Crefname{originallemma}{Original lemma}{Original lemmas}
\makeatletter
\newlabel{derived-lemma-pro-isomorphism-bis}{{0G3D}{source}{}{sourcepointer-2}{}}
\newlabel{derived-lemma-pro-isomorphism-bis@cref}{{[originallemma][0][]0G3D}{[1][0][]source}{}{}{}}
\makeatother
\crefname{originallemma}{original lemma}{original lemmas}
\Crefname{originallemma}{Original lemma}{Original lemmas}
\makeatletter
\newlabel{lemma-consequence-Artin-Rees-bis}{{0G3K}{source}{}{sourcepointer-3}{}}
\newlabel{lemma-consequence-Artin-Rees-bis@cref}{{[originallemma][0][]0G3K}{[1][0][]source}{}{}{}}
\makeatother
\title{The comparison between the derived tensor product with the powers of an ideal and the ordinary products is a pro-isomorphism}
\author{The Stacks Project Authors}
\date{stacks-project:more-algebra.tex:local-snapshot-20261012; GNU FDL 1.2 or later}
\begin{document}
\maketitle
\setcounter{subsection}{1}
\noindent\textit{Editorial note.} All delivered mathematical text is the author's own, copied from the pinned Stacks Project chapter file more-algebra.tex; no statement or proof step was rewritten, reordered or removed. The counted statement is the labeled lemma environment at source lines 29676-29685 (canonical tag 0G3N) and the counted proof is the proof environment at source lines 29687-29734, which immediately follows it and is its own proof.\par\medskip
\noindent\textit{Editorial note.} The proof contains one sentence that is not terminated by a period before the next display begins, at source line 29722; the printed text is copied exactly as it stands and is not corrected.\par\medskip
\noindent\textit{Editorial note.} The proof refers to lemma-sequence-powers-pro-bounded and lemma-consequence-Artin-Rees-bis of this chapter and to a pro-isomorphism criterion of another chapter; all three are carried as navigation stubs with their roles recorded in the reasons.\par\medskip
\noindent\textit{Editorial note.} The chapter file pinned-original-source.tex carries only an input of the shared preamble on its first line, so no preamble text lives inside the chapter. The retained file sources/stacks-more-algebra/retained-preamble.tex is a mechanical, line-for-line copy of the project's single shared preamble as staged in the snapshot (Xy-pic setup plus the native theorem declarations and mathematical macros). Nothing was added, removed or redefined, and the excerpt serves only as preamble, never as a span coordinate.\par\medskip
\noindent\textit{Editorial note.} The source URL is the source-level Stacks Project URL, as required when more than one card is produced from one source; the per-claim tag is carried by claim\_key and by this note. The counted statement and proof cite no bibliography entry, so this card needs no bibliography.\par\medskip
\noindent\textit{Allowed standard prerequisites (editorial).}\begin{itemize}
\item The inverse system of powers I\textasciicircum{}n is pro-isomorphic to the system generated by the n-th powers of a finite set of generators, and the auxiliary systems attached to the generators carry the same pro-information. used at proof lines 29688-29698. Replaces the original system by the system generated by the chosen generators and reduces the statement to that system.
\item An isomorphism of pro-objects of a derived category can be checked on cohomology in each degree, and it suffices to compute the derived tensor product with a bounded above complex of finite free modules in place of the bounded complex M. used at proof lines 29699-29725. Reduces the pro-isomorphism of complexes to a pro-isomorphism of modules in each cohomological degree, for a convenient representative of M.
\item The Artin-Rees consequence supplies the pro-isomorphism between the cohomology of the tensor product with the powers of the generators and the cohomology of the ordinary products. invoked at proof lines 29731-29733. Provides the final pro-isomorphism, thereby proving the lemma.
\end{itemize}\medskip
\ifdefined\hypertarget\hypertarget{sourcepointer-1}{}\fi
\noindent\textit{Source reference.} Original lemma 0G3M: sources/stacks-more-algebra/pinned-original-source.tex, source lines 29648--29648. The reference supplies the pro-isomorphism used at proof lines 29691-29692 to replace the system of powers by the system generated by f\_1, ..., f\_r; the reductions that use it are inside the counted span.\par\medskip
\ifdefined\hypertarget\hypertarget{sourcepointer-2}{}\fi
\noindent\textit{Source reference.} Original lemma 0G3D: sources/stacks-chow/tags.txt, source lines 20088--20088. The reference is the criterion that lets an isomorphism of pro-objects be checked on cohomology in each degree, invoked at proof lines 29709-29710; the degree bound that makes the check sufficient is established inside the counted span.\par\medskip
\ifdefined\hypertarget\hypertarget{sourcepointer-3}{}\fi
\noindent\textit{Source reference.} Original lemma 0G3K: sources/stacks-more-algebra/pinned-original-source.tex, source lines 29432--29432. The reference is the Artin-Rees consequence applied at proof lines 29731-29733 to the last comparison of cohomology; the three reductions that bring the lemma into the form where this consequence applies are inside the counted span.\par\medskip
\begin{lemma}
\label{lemma-tensoring-Deligne-system}
Let $A$ be a Noetherian ring. Let $I \subset A$ be an ideal. Let
$M^\bullet$ be a bounded complex of finite $A$-modules. The
inverse system of maps
$$
I^n \otimes_A^\mathbf{L} M^\bullet \longrightarrow I^nM^\bullet
$$
defines an isomorphism of pro-objects of $D(A)$.
\end{lemma}

\begin{proof}
Choose generators $f_1, \ldots, f_r \in I$ of $I$. The inverse system $I^n$
is pro-isomorphic to the inverse system $(f_1^n, \ldots, f_r^n)$
in the category of $A$-modules. With notation as in
Lemma \ref{lemma-sequence-powers-pro-bounded} we find that
it suffices to prove the inverse system of maps
$$
I_n^\bullet \otimes_A^\mathbf{L} M^\bullet
\longrightarrow
(f_1^n, \ldots, f_r^n)M^\bullet
$$
defines an isomorphism of pro-objects of $D(A)$.
Say we have $a \leq b$ such that $M^i = 0$ if $i \not \in [a, b]$.
Then source and target of the arrows above have cohomology
only in degrees $[-r + a, b]$.
Thus it suffices to show that for any $p \in \mathbf{Z}$ the
inverse system of maps
$$
H^p(I_n^\bullet \otimes_A^\mathbf{L} M^\bullet)
\longrightarrow
H^p((f_1^n, \ldots, f_r^n)M^\bullet)
$$
defines an isomorphism of pro-objects of $A$-modules, see
Derived Categories, Lemma \ref{derived-lemma-pro-isomorphism-bis}.
Using the pro-isomorphism between
$I_n^\bullet \otimes_A^\mathbf{L} M^\bullet$
and $I^n \otimes_A^\mathbf{L} M^\bullet$
and the pro-isomorphism between
$(f_1^n, \ldots, f_r^n)M^\bullet$ and $I^nM^\bullet$
this is equivalent to showing that the inverse system of maps
$$
H^p(I^n \otimes_A^\mathbf{L} M^\bullet)
\longrightarrow
H^p(I^nM^\bullet)
$$
defines an isomorphism of pro-objects of $A$-modules
Choose a bounded above complex of finite free $A$-modules
$P^\bullet$ and a quasi-isomorphism $P^\bullet \to M^\bullet$.
Then it suffices to show that the inverse system of maps
$$
H^p(I^nP^\bullet)
\longrightarrow
H^p(I^nM^\bullet)
$$
is a pro-isomorphism. This follows from
Lemma \ref{lemma-consequence-Artin-Rees-bis}
as $H^p(P^\bullet) = H^p(M^\bullet)$.
\end{proof}

\end{document}
